\documentclass[11pt]{article}
\usepackage{fullpage,url,hyperref}
\usepackage{amsmath}
\usepackage{amssymb}

\usepackage[letterpaper,top=1in,bottom=1in,left=1in,right=1in,nohead]{geometry}

\setlength{\parindent}{0in}
\setlength{\parskip}{6pt}

\begin{document}
\thispagestyle{empty}
{\large{\bf ECE 6701 + CS 6501: Optimization for Engineering \hfill Due Fri 9/18}}\\

{\LARGE{\bf Homework 1: Convex Sets and Functions}}
\vspace{0.2\baselineskip}
\hrule

{\bf Instructions:} Submit a single Jupyter notebook (.ipynb) of your work to
Canvas by 11:59pm on the due date. All math must be formatted using LaTeX. {\bf Be sure to show all the work involved in deriving your answers! If you just give
  a final answer without explanation, you may not receive credit for that
  question.}

You may discuss the concepts with your classmates, look on the internet for
help, or ask an LLM for advice, but first try hard to solve problems on your
own, and always write up the answers entirely on your own. Make sure you
understand the concepts and solutions---don't just look up the answers or use AI
to answer it for you. {\bf Cite any sources you used outside of the class
  material (webpages, LLMs, etc.).}

\begin{enumerate}

\item {\bf Dual Norms}
Let $\| \cdot \| : \mathbb{R}^n \rightarrow \mathbb{R}$ be a norm. Define the dual norm, denoted $\| \cdot \|_*$, as follows for any vector $z \in \mathbb{R}^n$:
$$\|z\|_* = \sup \{z^T x : \|x\| \leq 1 \}.$$

\begin{enumerate}
  \item Show that $\| \cdot \|_*$ is a norm, i.e., show it satisfies the four requirements to be a norm that we covered in class.

{\bf Hint:} You may use the following fact (without proving it yourself). Let $S$ be a set and $f : S \rightarrow \mathbb{R}$, $g : S \rightarrow \mathbb{R}$. Then,
$$\sup \{ f(s) + g(s) : s \in S\} \leq \sup \{f(s) : s \in S\} + \sup \{g(s) : s \in S\}$$

\item Show that the dual of the dual of a norm is the original norm. That is, for any norm $\| \cdot \|$, we have
$$\|x\|_{**} = \|x\|, \quad \text{for all $x \in \mathbb{R}^n$}.$$

\item Show that the dual of the $L^2$ norm (a.k.a. Euclidean norm) is again the $L^2$ norm. That is,
$$\|x\|_{2*} = \|x\|_2, \quad \text{for all $x \in \mathbb{R}^n$}.$$
{\bf Hint:} You might find the Cauchy-Schwarz inequality useful!

\end{enumerate}

{\bf Not required---don't turn in!} If you want to go further, you might
  be interested to prove to yourself the following facts (or look up the
  proofs):
  \begin{itemize}
  \item The dual of the $L^1$ norm is the $L^\infty$ norm.
  \item The dual of the $L^p$ norm is the $L^q$ norm, where
    $\frac{1}{p} + \frac{1}{q} = 1$. This follows from a relation known as
    Holder's Inequality, which is a generalization of the Cauchy-Schwarz
    Inequality.  See Appendix A.1.6 (dual norms) and Section 3.1.9 (Holder's
    Inequality) of the B\&V book and also the Wikipedia pages on ``Dual Norm,''
    ``$L^p$ Space,'' and ``Holder's Inequality'' for more information.
  \end{itemize}
  

\item {\bf Point on a Hyperplane Closest to the Origin}
  
  Let $a \in \mathbb{R}^n, a \neq 0$, and $b \in \mathbb{R}$. Define the hyperplane
  $$S = \{x \in \mathbb{R}^n : a^T x = b\}.$$
  Give an expression for the point $x_0 \in S$ that is closest to the origin in Euclidean distance, or said another way, the point where the $L^2$ norm is minimal:
  $$x_0 = \arg \min_{x \in S} \|x\|_2.$$ {\bf Hint:} The formula for $x_0$ should be an expression in terms of $a$ and $b.$

\item {\bf Intersection of Hyperplanes}

  Consider two hyperplanes, defined by $a_1, a_2 \in \mathbb{R}^n, a_1 \neq 0, a_2 \neq 0,$ and $b_1, b_2 \in \mathbb{R}:$
  \begin{align*}
    S_1 &= \{ x \in \mathbb{R}^n : a_1^T x = b_1\}, \\
    S_2 &= \{ x \in \mathbb{R}^n : a_2^T x = b_2\}.
  \end{align*}

  \begin{enumerate}
    \item Under what conditions are these hyperplanes disjoint, i.e., $S_1 \cap S_2 = \emptyset$?
      (State the conditions precisely in terms of the $a_i$ and $b_i$.)

    \item When $S_1$ and $S_2$ do intersect, what kind of space is $S_1 \cap S_2$?
      What is its dimension? Please list all possible cases ({\bf hint:} there are two
      cases, with two different dimensions), and describe the conditions under which
      each case happens.

\end{enumerate}


\item {\bf Convex Cones from Convex Sets}
  
  Let $S \subset \mathbb{R}^n$ be a convex set. Define the set
  $$C = \{\theta x : x \in S, \theta \geq 0\}.$$
  Show that $C$ is a convex cone.

\item {\bf Convex Combinations and Convex Functions}

  \begin{enumerate}

  \item Let $C \subset \mathbb{R}^n$ be a convex set, and let $f : C \rightarrow \mathbb{R}$
    be a convex function. Show that for any $x_1, \ldots, x_k \in C$ and any
    $\theta_1, \ldots, \theta_k \geq 0$, with $\theta_1 + \cdots + \theta_k = 1$, 
    $$f(\theta_1 x_1 + \cdots \theta_k x_k) \leq \theta_1 f(x_1) + \cdots + \theta_k f(x_k).$$

  \item Let $f : \mathbb{R}^n \rightarrow \mathbb{R}$ be a convex function. Show that for any
    $x_1, \ldots, x_k \in \mathbb{R}^n$ and any $\theta_1 > 0,$ and $\theta_2, \ldots, \theta_k \leq 0$,
    with $\theta_1 + \cdots + \theta_k = 1,$
    $$f(\theta_1 x_1 + \cdots + \theta_k x_k) \geq \theta_1 f(x_1) + \cdots + \theta_k f(x_k).$$
    ({\bf Careful:} Pay close attention to the direction of the various inequality signs!)

\end{enumerate}

\item {\bf Convex Function or Not?}

For each of the following functions, $f : \mathbb{R}^n \rightarrow \mathbb{R}$,
state whether it is convex, concave, both, or neither. Justify your answer in
each case.
\begin{enumerate}
\item $f(x) = x^T x - \frac{1}{n} (1^T x)^2$
\item $f(x) = \|x\|_1 + \frac{1}{2} \|x\|_2^2$
\item
$f(x) = \log \frac{\exp(a^T x + b)}{1 + \exp(a^T x + b)}, \quad \text{for $a \in \mathbb{R}^n \backslash \{0\}, b \in \mathbb{R}$}$
\end{enumerate}

\item {\bf Epigraphs of Functions}

Problem 3.6 from B\&V: When is the epigraph of a function a halfspace? When is the
epigraph of a function a convex cone? When is the epigraph of a function a polyhedron?

\end{enumerate}

\section*{Bonus \LaTeX Example}
Here is a simple proof of the Cauchy-Schwarz inequality, just for an example of
how to use the {\tt align} environment to do a step-by-step calculation.

Given two vectors $v = (v_1, \ldots, v_n) \in \mathbb{R}^n$ and $w = (w_1,
\ldots, w_n) \in \mathbb{R}^n$, our goal is to show the Cauchy-Schwarz
inequality:
$$|\langle v, w \rangle| \leq \|v\|_2 \|w\|_2.$$
We start with rewriting the following quadratic:
\begin{align*}
  0 & \leq \sum_i \sum_j (v_i w_j - v_j w_i)^2 & \text{non-negative because it's a square}\\
    & = \sum_i \sum_j v_i^2 w_j^2 - 2 v_i w_i v_j w_j + v_j^2 w_i^2 & \text{expanding terms}\\
    & = 2 \|v\|_2^2 \|w\|_2^2 - 2\langle v, w \rangle^2 & \text{writing sums as inner product / norm}\\
\end{align*}
Now, by moving the last term to the other side of the inequality, we get:
$$\langle v, w \rangle^2 \leq \|v\|_2^2 \|w\|_2^2.$$
Taking the square-root of both sides gives us the Cauchy-Schwarz inequality.
\end{document}
